
\chapter{Conclusion}
This thesis has explored how exposure to international maritime trade shapes the likelihood of conflict across Nigerian regions. The analysis combined information on port-level activity, maritime connectivity, and geography to build an index of local exposure to global trade shocks. The evidence points to a heterogeneous pattern: the effect of trade on conflict varies widely across space, sectors, and local resource conditions.

\section{Main findings}
\textit{Baseline and falsification.} For commercial ports, the aggregate effect of the exposure index is small on average, but analysing the heterogeneity of the sample given the development of each ADM2 (proxied by night-lights), we discovered that in relatively less developed areas, higher exposure to trade is associated with lower conflict, while effects in more developed areas are near zero. When index is extentet to all ports (including those without commercial relevance), coefficients become weaker and statistically insignificant. This falsification supports the interpretation that the mechanism runs through international trade, in particolar access to international market, rather than mere proximity to coastal infrastructure. \\

\textit{Geographic heterogeneity.} We studied different heterogeneity to study how the relationship varies across different dimension. Results show that there is a stabilizing effect stronger in the eastern belt than in the west. This spatial asymmetry matches economic patterns: eastern ADM2s are more rural, as shown by the proxy of nightlights, making them more sensitive to trade-induced welfare improvements.\\

\textit{Sector and trade flow heterogeneity.} Result shows that the negative association is more pronounced for ports oriented toward the one that trade vegetables, as second most traded industry, and imports more with respect of what they export. This is consistent with our theory of consumption-smoothing channel.\\

\textit{Resource endowments (water).} Interactions with water availability (a proxy for agricultural potential) indicate that the pacifying effect concentrates in resource-constrained areas. Where water availability is higher—hence greater agricultural rents or competition—the relationship can weaken or even turn positive, suggesting that trade may intensify distributional tensions when valuable rents are at stake. \\


\section{Interpretation}
Overall, the findings support an opportunity-income channel, especially in poorer districts. Where communities gain easier access to markets or cheaper imported goods, the economic incentives to engage in violence tend to weaken. In contrast, where economic rents and pre-existing advantages are larger, trade may amplify inequalities and make consumption-related conflicts more likely. In short, trade is stabilising where markets fill and income opportunities, while it risks destabilising where it amplifies returns and contestable surpluses. 

\section{Limitations}
While the findings are robust across specifications reported here, several caution should be taking in consideration:
First measurement. CPT and LSCI are port indicators available for a shorter window, for which we mapped to the study period and the other port indicators, sectoral and flow classifications rely on IMF PortWatch assumption. Night-lights and water availability are proxies for income and agricultural potential rather than direct measures.
Second, inverse distance to the nearest port abstracts from the situation of actual infrastructures; future work should embed explicit transport networks and travel times.
Third conflict counts aggregate heterogeneous events; although the interpretation emphasizes consumption-related mechanisms, misclassification remains possible.
Lastely, the design mitigates simultaneity by using lagged port shocks and fixed effects, and the “all-ports” falsification helps; still, unobserved time-varying shocks correlated with both exposure and conflict cannot be fully ruled out.

\section{Policy implications}
The heterogeneity of effects suggests that how a country integrates is as important as the volume of its trade.
\textit{Target lagging regions}. Investments that reduce effective distance to commercial ports for poorer inland districts can enhance the conflict-reducing benefits of trade.\\
\textit{ Safeguard the welfare channel}. Policies that keep essential import pipelines reliable and limit pass-through frictions help stabilize local prices and incomes.\\
textit{Manage rents where activity is high}. In more developed or water-abundant areas, complement trade expansion with transparency, competition policy, and inclusive local fiscal instruments to curb rent capture and distributional conflict.
\textit{Shock-response tools}. Temporary safety nets and price-stabilization measures during large external swings can preserve the pacifying effect of integration.\\

\section{External validity and future research}
Although Nigeria’s coastal geography and regional fragmentation make it a compelling test case, external validity should be assessed cautiously. Extending the approach to other coastal African economies and to land-linked settings connected by corridors would clarify the role of transport networks. Three concrete extensions emerge: (i) network-based exposure, in fact replace straight-line distance with multimodal travel times and port-choice models; allow exposure to multiple ports and transshipment hubs. (ii) Micro-outcomes such as link exposure to prices, firm dynamics, and employment to pin down the welfare more directly. (iii) Institutional moderators: interact exposure with local governance quality, security force presence, and customary authority to test capacity-to-manage-distributional-tensions.

\section{Contribution}
Empirically, the thesis introduces a geography aware exposure metric that ties port level shocks and connectivity to local conflict dynamics, and documents sizable heterogeneity by development, location, sector, and resource endowment. Conceptually, it bridges the literature on external income shocks and the spatial diffusion of global markets: trade can pacify—but not everywhere and not unconditionally. Where integration expands livelihoods, conflict tends to recede; where it magnifies contestable rents, stability may not follow. Recognizing and acting on this conditionality is central to leveraging globalization for peace.

